% Convertido de estado_da_arte.md.
AAUTIA --- Aprendizagem Automática \textbar{}
PLANTIA --- Planeamento em Inteligência Artificial

\textbf{Tema do projeto:} \emph{Sistema Inteligente de Apoio à Decisão
de Treinadores de Basquetebol através de Machine Learning e Planeamento
Automático.}

Esta primeira versão do estado da arte apresenta \textbf{trabalhos científicos reais, publicados e com referências verificáveis}, sobre previsão de resultados, avaliação de jogadores, composição de equipas, substituições, otimização e planeamento de estratégias.

A investigação foi organizada em duas partes, seguindo a divisão das unidades curriculares:

\begin{itemize}
\tightlist
\item
  \textbf{AAUTIA (50\%)} --- modelos preditivos, aprendizagem   supervisionada, avaliação de desempenho e previsão de resultados.
\item
  \textbf{PLANTIA (50\%)} --- escolha de alinhamentos, substituições,   otimização de decisões, processos de Markov e planeamento sob   incerteza.
\end{itemize}

Os trabalhos foram identificados através de editoras científicas, revistas académicas, repositórios universitários e publicações com DOI. 
A apresentação é feita em português, mantendo os títulos originais para facilitar a localização.

\section{Estado da arte de AAUTIA}\label{aautia:estado-da-arte-de-aautia}

\subsection{Previsão de resultados no basquetebol
europeu}\label{aautia:previsuxe3o-de-resultados-no-basquetebol-europeu}

Lampis, T., Ntzoufras, I., Vassalos, V. e Dimitriou, S. (2023)
\cite{lampis2023}

\textbf{Artigo:} \emph{Predictions of European Basketball Match Results with Machine Learning Algorithms}

\textbf{Publicação:} Journal of Sports Analytics.

\href{https://doi.org/10.3233/JSA-220639}{Consultar artigo científico --- DOI}

\subsubsection{Objetivo e metodologia}\label{aautia:objetivo-e-metodologia}

Este trabalho investiga a previsão de resultados de jogos de basquetebol em quatro competições europeias: EuroLeague, EuroCup, Liga Grega e Liga Espanhola ACB.

Os investigadores utilizaram \textbf{5 214 jogos realizados entre 2013 e 2018}, construindo indicadores a partir das estatísticas históricas das equipas.

Foram comparados vários modelos, nomeadamente:

\begin{itemize}
\tightlist
\item
  Regressão Logística;
\item
  Random Forest;
\item
  XGBoost;
\item
  Combinações de modelos (\emph{ensemble learning}).
\end{itemize}

O estudo utilizou ainda sistemas de classificação de equipas, como Elo, PageRank e \emph{pi-rating}, e indicadores de desempenho para enriquecer os dados de entrada.

\subsubsection{Resultados obtidos}\label{aautia:resultados-obtidos}

Os modelos com indicadores adicionais superaram os modelos de referência mais simples em aproximadamente 3 a 5 pontos percentuais de exatidão.

Não foi identificado um algoritmo universalmente superior. O desempenho dependeu da competição e das características dos dados.

Os autores também verificaram a importância dos indicadores de forma recente e dos sistemas de classificação na previsão dos resultados.

\subsubsection{Relevância para AAUTIA}\label{aautia:relevuxe2ncia-para-aautia}

Este é um dos trabalhos mais importantes para o nosso projeto porque fornece uma metodologia diretamente reproduzível: construir características históricas, comparar algoritmos e avaliar previsões com métricas como \emph{accuracy}, \emph{F1-score} e \emph{Brier score}.

\textbf{Aplicação proposta:} utilizar Regressão Logística, Random Forest e XGBoost para prever vitórias, comparando modelos baseados apenas em resultados com outros que incluam estatísticas das equipas e indicadores individuais.

\subsection{Integração de XGBoost e IA explicável na previsão de resultados da NBA}\label{aautia:integrauxe7uxe3o-de-xgboost-e-ia-explicuxe1vel-na-previsuxe3o-de-resultados-da-nba}

Ouyang, Y., Li, X., Zhou, W., Hong, W., Zheng, W., Qi, F. e Peng, L. (2024) \cite{ouyang2024}

\textbf{Artigo:} \emph{Integration of Machine Learning XGBoost and SHAP Models for NBA Game Outcome Prediction and Quantitative Analysis Methodology}

\textbf{Publicação:} PLOS ONE, 19(7), e0307478.

\href{https://doi.org/10.1371/journal.pone.0307478}{Consultar artigo completo --- acesso aberto}

\subsubsection{Objetivo e metodologia}\label{aautia:objetivo-e-metodologia-1}

Os autores desenvolveram um sistema de previsão de resultados da NBA utilizando dados das temporadas de 2021 a 2023.

O trabalho combina duas tecnologias:

\textbf{XGBoost:} algoritmo de aprendizagem supervisionada baseado em conjuntos de árvores de decisão, utilizado para prever resultados.

\textbf{SHAP:} técnica de explicabilidade que permite estimar a contribuição das diferentes variáveis para as previsões realizadas pelo modelo.

O sistema analisa indicadores de desempenho e permite atualizar as previsões em diferentes momentos de uma partida.

\subsubsection{Resultados e conclusões}\label{aautia:resultados-e-conclusuxf5es}

O estudo demonstra uma abordagem para prever resultados e identificar quais os indicadores estatísticos que mais influenciam as previsões.

A integração de SHAP permite interpretar os modelos, em vez de apresentar apenas uma probabilidade de vitória.

Os autores disponibilizaram também o código utilizado na investigação.

\href{https://github.com/YanOuyang514/NBA-game-outcome-prediction-and-quantitative-analysis-methodology}{Código do trabalho no GitHub}

\subsubsection{Relevância para AAUTIA}\label{aautia:relevuxe2ncia-para-aautia-1}

Este artigo é especialmente útil porque o nosso projeto não deve limitar-se a informar o treinador de que existe determinada probabilidade de vitória.

Deve explicar \textbf{quais os fatores que contribuem para essa previsão}.

Por exemplo, o sistema poderá indicar que uma previsão foi influenciada pelo desempenho defensivo recente, pela eficiência dos lançamentos e pelo rendimento histórico de determinados alinhamentos.

Importa distinguir estas explicações estatísticas de afirmações causais sobre o que acontecerá se o treinador alterar uma decisão.

\subsection{Previsão dos resultados da NBA com Machine Learning}\label{aautia:previsuxe3o-dos-resultados-da-nba-com-machine-learning}

Horvat, T., Job, J., Logozar, R. e Livada, Č. (2023) \cite{horvat2023}

\textbf{Artigo:} \emph{A Data-Driven Machine Learning Algorithm for Predicting the Outcomes of NBA Games}

\textbf{Publicação:} Symmetry, 15(4), 798.

\href{https://doi.org/10.3390/sym15040798}{Consultar artigo científico}

\subsubsection{Objetivo e metodologia}\label{aautia:objetivo-e-metodologia-2}

Este trabalho propõe um modelo baseado em dados para prever resultados de jogos da NBA.

Os autores dão particular importância à construção de indicadores estatísticos, procurando transformar os dados disponíveis em variáveis com capacidade preditiva.

O estudo estabelece uma formulação matemática para as estatísticas de basquetebol e para os indicadores derivados, utilizando-os na construção de modelos de aprendizagem automática.

\subsubsection{Contributos}\label{aautia:contributos}

Um dos aspetos relevantes é a preocupação com a definição rigorosa das variáveis utilizadas nos modelos.

A qualidade das características extraídas dos dados é determinante para a capacidade de generalização das previsões.

Este trabalho reforça a importância de construir uma base de dados coerente antes de escolher algoritmos mais complexos.

\subsubsection{Relevância para AAUTIA}\label{aautia:relevuxe2ncia-para-aautia-2}

No nosso projeto, poderá servir de referência para construir indicadores como:

\begin{itemize}
\tightlist
\item
  eficiência ofensiva e defensiva;
\item
  média de pontos nos últimos jogos;
\item
  aproveitamento dos lançamentos;
\item
  ressaltos ofensivos e defensivos;
\item
  assistências e perdas de bola.
\end{itemize}

Estes indicadores devem ser calculados apenas com informação disponível antes de cada partida, evitando a utilização inadvertida de dados futuros.

\subsection{Previsão de resultados durante o decorrer do jogo}\label{aautia:previsuxe3o-de-resultados-durante-o-decorrer-do-jogo}

Kandhway, K. (2025) \cite{kandhway2025}

\textbf{Artigo:} \emph{In-Play Outcome Prediction of an Ongoing Basketball Game Using Machine Learning}

\textbf{Publicação:} Journal of Prediction Markets, 18(2).

\href{https://doi.org/10.5750/jpm.v18i2.2155}{Consultar o artigo científico}

\subsubsection{Objetivo e metodologia}\label{aautia:objetivo-e-metodologia-3}

Este trabalho investiga a previsão do vencedor de um jogo de basquetebol \textbf{enquanto a partida está a decorrer}.

Ao contrário dos modelos pré-jogo, o sistema atualiza as previsões com base na evolução da partida.

Foram utilizados dados da NBA correspondentes a oito temporadas, de 2015/2016 a 2022/2023.

O estudo integra técnicas de aprendizagem supervisionada e utiliza SHAP (SHapley Additive exPlanations) para interpretar a importância das diferentes variáveis.

\subsubsection{Resultados obtidos}\label{aautia:resultados-obtidos-1}

Os autores reportam uma exatidão global de aproximadamente \textbf{74\%}, com resultados que variam entre cerca de 63\% e 90\% em diferentes estados do jogo.

Os resultados mostram também que a importância das variáveis se altera ao longo da partida.

\subsubsection{Relevância para AAUTIA}\label{aautia:relevuxe2ncia-para-aautia-3}

Este trabalho fundamenta uma evolução do nosso sistema para decisões em tempo real.

Por exemplo, o treinador poderia receber uma previsão atualizada após o terceiro período, considerando a diferença pontual, o tempo restante e as estatísticas acumuladas.

Numa primeira fase, considero preferível manter a previsão pré-jogo; a previsão dinâmica poderá constituir uma experiência complementar.

\subsection{Revisão sistemática de aplicações de IA à previsão de basquetebol}\label{aautia:revisuxe3o-sistemuxe1tica-de-aplicauxe7uxf5es-de-ia-uxe0-previsuxe3o-de-basquetebol}

Li, C., Zhang, H., Zhang, Y., Shen, J. e An, R. (2025) \cite{li2025}

\textbf{Artigo:} \emph{The Application of Artificial Intelligence Techniques in Predicting Game Outcomes of Professional Basketball League: A Systematic Review}

\textbf{Publicação:} PLOS ONE, 20(6), e0326326.

\href{https://doi.org/10.1371/journal.pone.0326326}{Consultar revisão sistemática completa}

\subsubsection{Objetivo e metodologia}\label{aautia:objetivo-e-metodologia-4}

Esta revisão analisa a aplicação de métodos de Inteligência Artificial à previsão de resultados em ligas profissionais de basquetebol.

Foram incluídos \textbf{34 estudos científicos}, com uma janela de pesquisa centrada em trabalhos publicados entre 2019 e 2024.

A revisão compara diferentes abordagens, incluindo redes neuronais, Random Forest e modelos combinados.

\subsubsection{Conclusões}\label{aautia:conclusuxf5es}

Os autores identificam resultados promissores para diferentes algoritmos e salientam a importância da seleção de características e da preparação dos dados.

Recomendam também melhorias na validação dos modelos, maior diversidade de competições analisadas e investigação sobre previsão dinâmica e transferência de aprendizagem.

Alguns trabalhos analisados reportam valores de exatidão muito elevados. Porém, esses resultados não devem ser entendidos como desempenhos garantidos ou diretamente comparáveis, porque os conjuntos de dados, as variáveis e os procedimentos de validação diferem.

\subsubsection{Relevância para AAUTIA}\label{aautia:relevuxe2ncia-para-aautia-4}

Este será um artigo central da introdução e do enquadramento científico.

Permite justificar a seleção dos algoritmos, identificar abordagens já exploradas e estabelecer critérios de avaliação para o trabalho experimental.
